\chapter*{Come usare questo formulario}
\addcontentsline{toc}{chapter}{Come usare questo formulario}

\begin{obiettivocapitolo}{risolvere, non soltanto ricordare}
Ogni esercizio va affrontato in quattro passaggi: \textbf{riconoscimento}, \textbf{scelta}, \textbf{esecuzione}, \textbf{controllo}. Le formule sono strumenti; non sostituiscono la verifica del dominio, delle ipotesi e dei casi esclusi.
\end{obiettivocapitolo}

\begin{teoriaessenziale}{confine del volume}
Le equazioni differenziali ordinarie (EDO) non sono trattate in questo formulario. Anche quando compaiono in alcuni programmi universitari alla fine di Analisi I, nella collana di ingegnerismo.it appartengono al \href{https://ingegnerismo.it/matematica/formulario-edo-analisi-numerica/}{Formulario di EDO e Analisi Numerica}, così da evitare duplicazioni e sovrapposizioni tra volumi.
\end{teoriaessenziale}

\begin{proceduraesame}{sequenza universale in quattro fasi}
\begin{enumerate}
\item \textbf{Riconosci la famiglia del problema.} Individua oggetto, dati, incognita, dominio e regime: equazione, limite, derivata, integrale, serie, approssimazione oppure studio globale.
\item \textbf{Scegli il metodo.} Usa le tabelle decisionali all'inizio di ciascun capitolo. Se un test è inconcludente, passa al criterio successivo invece di forzare una conclusione.
\item \textbf{Esegui in forma controllabile.} Scrivi le trasformazioni importanti, le ipotesi utilizzate e i valori esclusi. Mantieni le espressioni esatte fino all'ultimo passaggio, salvo richiesta numerica.
\item \textbf{Controlla.} Verifica dominio, segno, ordine di grandezza, eventuali rami, condizioni iniziali, derivata della primitiva, comportamento agli estremi o convergenza del criterio applicato.
\end{enumerate}
\end{proceduraesame}

\begin{sceltametodo}{come leggere i riquadri}
\begin{tabularx}{\textwidth}{@{}>{\bfseries\sffamily}p{.23\textwidth}X@{}}
Scelta del metodo & riconosce la forma della traccia e indica il primo criterio da tentare.\\
Procedura d'esame & fornisce passaggi ordinati, ripetibili e verificabili.\\
Formula chiave & raccoglie la relazione da usare con le condizioni essenziali.\\
Errore frequente & segnala passaggi formalmente plausibili ma matematicamente illegittimi.\\
Controllo finale & indica verifiche rapide prima di consegnare.
\end{tabularx}
\end{sceltametodo}

\begin{errorefrequente}{saltare direttamente alla formula}
Prima di sostituire numeri o manipolare simboli, controllare sempre: denominatori non nulli, argomenti dei logaritmi positivi, radicandi delle radici pari non negativi, intervallo di definizione, orientazione degli integrali e ipotesi dei teoremi.
\end{errorefrequente}
